\chapter{Urn Model}
\label{ch:3}
\section{Markov chains}
\label{sec:3-1}
In the discrete-time Ehrenfest chain, the state $j$ is the number of balls in one of two urns. Moving one uniformly chosen ball gives
\[
(Pf)(j)=\frac{N-j}{N}f(j+1)+\frac jN f(j-1).
\]
The Krawtchouk difference equation gives
\[
P K_n(\cdot;N)=\left(1-\frac{2n}{N}\right)K_n(\cdot;N).
\]
Here $N\geq1$, and terms outside the state space have zero transition coefficient.

\section{Initial and invariant distributions}
\label{sec:3-2}
The invariant distribution is $w_N(j)=2^{-N}\binom Nj$. Detailed balance follows from
\[
w_N(j)\frac{N-j}{N}=w_N(j+1)\frac{j+1}{N}.
\]
The non-lazy discrete-time chain has period two; invariance should not be confused with convergence to equilibrium from every initial distribution.



\section{Induced matrices and multinomial chains: the SFA II framework}
\label{sec:ency-sfa-ii}
The multinomial extension of a one-site chain is the induced-matrix
construction of SFA II \cite{FeinsilverSFAII2023}. It provides the finite
framework for the cumulative-trial example below and the Mehler limit
in Section~\ref{sec:11-3}. We declare row-stochastic conventions here;
column-vector probability evolution uses the transpose.

Use $u,v$ for category indices, reserving $a,b$ for success counts below.
For a $q$-state matrix $A$, its degree-$N$ induced matrix is defined by
\begin{equation}\label{eq:ency-induced-definition}
\prod_u\left(\sum_vA_{uv}z_v\right)^{n_u}
=\sum_{|\mathbf m|=N}A^{[N]}(\mathbf n,\mathbf m)\mathbf z^{\mathbf m},
\qquad |\mathbf n|=N.
\end{equation}
Substitution gives $(AB)^{[N]}=A^{[N]}B^{[N]}$. If $A=P$ is stochastic,
$P^{[N]}$ is the transition for $N$ independent sites with their labels
forgotten. Equivalently,
\[
P^{[N]}(\mathbf n,\mathbf m)
=\sum_{\substack{\sum_v c_{uv}=n_u\\\sum_u c_{uv}=m_v}}
\prod_u\frac{n_u!}{\prod_v c_{uv}!}\prod_{u,v}P_{uv}^{c_{uv}}.
\]
Thus the contingency-table model and the induced matrix are exactly
the same finite construction. The coordinate matrix of the symmetric
tensor power depends on basis normalization; the stochastic occupation
basis should not be confused with an orthonormal tensor basis.

If $\boldsymbol\eta$ is stationary, then
\[
\Pi_N(\mathbf n)=\frac{N!}{\prod_u n_u!}\prod_u\eta_u^{n_u}
\]
is stationary. For a complete one-site eigenbasis with eigenvalues
$1,r_1,\ldots,r_{q-1}$, symmetric products give occupation eigenvalues
$\prod_{j=1}^{q-1}r_j^{\nu_j}$ for $|\boldsymbol\nu|\le N$.
Reversibility permits orthogonal one-site modes and the associated
multivariate Krawtchouk eigenbasis; stochasticity alone does not imply
this orthogonality.

For $P_t=e^{tQ}$ the induced generator is
\[
(\mathcal L_Nf)(\mathbf n)=\sum_{u\ne v}n_uQ_{uv}
[f(\mathbf n-e_u+e_v)-f(\mathbf n)].
\]
Cumulative Bernoulli trials specify a particular one-site $P$ and
therefore a particular $P^{[N]}$; they are not a different multinomial
extension. Their sampling interpretation makes positivity transparent,
while the induced construction explains composition and spectrum.
Matrix composition, polynomial linearization, and positive radial
convolution are distinct products, related only through their specified
spectral realizations.

Under stationary input the joint site probabilities are
$J_{uv}=\eta_uP_{uv}$. Recording $N$ such pairs gives a multinomial table
with margins $\mathbf N,\mathbf M$, and
\[
\frac{P^{[N]}(\mathbf n,\mathbf m)}{\Pi_N(\mathbf m)}
=\frac{\Pr(\mathbf N=\mathbf n,\mathbf M=\mathbf m)}
{\Pi_N(\mathbf n)\Pi_N(\mathbf m)}.
\]
This density ratio is the finite kernel whose central local limit is
Mehler's kernel. In the reversible case, diagonal fluctuation modes
give $q-1$ Gaussian factors and, after conjugation, $q-1$ oscillator
number operators. These limiting assertions require the local-limit
and mode-contraction arguments, beyond the induced-matrix identity.

\section{Cumulative Bernoulli trials: saved and recouped successes}
\label{sec:cumulative-bernoulli}
Hoare and Rahman \cite{HoareRahman1984} develop Krawtchouk processes
from combinations of Bernoulli trials with saved and recouped successes.
The earlier dynamic-occupancy program of Cooper, Hoare, and Rahman
\cite{CooperHoareRahman1977} starts from a Hahn kernel with
negative-hypergeometric factors. The following independent-site
construction is the binary specialization of Section~\ref{sec:ency-sfa-ii}
and presents our Bernoulli case with explicit conventions;
its parameters are defined here rather than implicitly identified with
all conventions in those papers.

Let $a$ and $b$ be the initial and final numbers of successes ($+1$ sites).
The corresponding negative-sign counts are $N-a$ and $N-b$; neither
success count is the negative-sign index $j$ of the defining Kravchuk
generating function. Starting from $a$ successes, save each success with
probability $\alpha$. If $k$ are saved, retry the remaining $N-k$ sites
with success probability $\beta$. With
$\operatorname{BinMass}(k;m,\theta)=\binom mk\theta^k(1-\theta)^{m-k}$ (zero outside $0\le k\le m$), the
positive transition kernel is
\begin{equation}\label{eq:ency-cumulative}
\mathcal K(b\mid a)=\sum_{k=0}^{\min(a,b)}
\operatorname{BinMass}(k;a,\alpha)\operatorname{BinMass}(b-k;N-k,\beta).
\end{equation}
Its backward action is $(\mathcal Kf)(a)=\sum_b\mathcal K(b\mid a)f(b)$;
the column-vector forward transition matrix is its transpose.
For $0\le\alpha<1$, $0<\beta<1$, the one-site transition and its
stationary parameter are
\[
P=\begin{pmatrix}1-\beta&\beta\\
(1-\alpha)(1-\beta)&\alpha+(1-\alpha)\beta\end{pmatrix},\qquad
r=\alpha(1-\beta),\quad p=\frac{\beta}{1-r}.
\]
Indeed $\mathbb E[Y-p\mid X]=r(X-p)$. The invariant occupation law is
$w_{N,p}(a)=\binom Na p^a(1-p)^{N-a}$, and, with $q=1-p$, the biased polynomials
$Q_n^{(p)}$ of \eqref{eq:biased-generating} in Section~\ref{sec:13-2}
satisfy
\[
\mathcal KQ_n^{(p)}=r^nQ_n^{(p)},\qquad
\sum_{a=0}^N w_{N,p}(a)Q_n^{(p)}(a;N)Q_m^{(p)}(a;N)=\binom Nn(pq)^n\delta_{nm}.
\]
Here the reference to $Q_n^{(p)}$ uses the current definition of
Section~\ref{sec:13-2}, evaluated at a success count; it does not identify
that count with the negative-sign index. Chapter~\ref{ch:13} is left
unchanged pending review. With its present definition, the explicit
coordinate bridge is
\[
Q_n^{(p)}(a;N)=(-1)^nQ_n^{(q)}(N-a;N).
\]
Both statements follow by expanding products of centered site variables
$X_u-p$. Thus the positive formula \eqref{eq:ency-cumulative} has the
exact bilinear expansion
\begin{equation}\label{eq:ency-cumulative-spectral}
\frac{\mathcal K(b\mid a)}{w_{N,p}(b)}
=\sum_{n=0}^N\frac{r^nQ_n^{(p)}(a;N)Q_n^{(p)}(b;N)}{\binom Nn(pq)^n}.
\end{equation}

For our symmetric convention, take $p=1/2$ and
$\beta=(1-r)/2$, $\alpha=2r/(1+r)$, with $0\le r<1$.
Then $P=\frac12\left(\begin{smallmatrix}1+r&1-r\\1-r&1+r\end{smallmatrix}\right)$
and $Q_n^{(1/2)}(a;N)=2^{-n}K_n(a;N)=(-1)^n2^{-n}K_n(N-a;N)$. For $r=e^{-t}$ this is the time-$t$ independent-bit
transition. It is not the single-flip Ehrenfest step: the spectra are
$r^n$ and $1-2n/N$, respectively.

At stationarity the $N$ input--output site pairs give a four-category
multinomial table. Its input and output margins $\mathsf S_{\mathrm{in}},\mathsf S_{\mathrm{out}}$ obey
\[
\frac{\mathcal K(b\mid a)}{w_{N,p}(b)}
=\frac{\Pr(\mathsf S_{\mathrm{in}}=a,\mathsf S_{\mathrm{out}}=b)}{w_{N,p}(a)w_{N,p}(b)}.
\]
This is the symmetric-tensor occupation kernel, interpreted as a joint
law relative to the product of its margins. Conditioning instead on a
fixed flip count gives hypergeometric overlap laws. Central Gaussian
scaling of the joint density ratio produces the Mehler kernel in
Section~\ref{sec:11-3}.

\section*{Connection to the supplied article sources}
Section~4 (``Ehrenfest urn model'') of the supplied
\texttt{KravFeinsKocik.tex} belongs to this chapter.
The urn state is the Hamming weight of a binary configuration: flipping
a uniformly selected coordinate moves from $j$ to $j+1$ with probability
$(N-j)/N$ and to $j-1$ with probability $j/N$, for $N\geq1$.
In the book's column-vector convention the transition matrix is
$A_N/N$. The tensor-space realization is discussed in Chapter~\ref{ch:9}.
